% Generated draft: requires mathematical and source review.
\documentclass[11pt]{article}
\usepackage[margin=25mm]{geometry}
\usepackage{fontspec}
\setmainfont{Latin Modern Roman}
\usepackage{amsmath,amssymb,mathtools,unicode-math}
\setmathfont{Latin Modern Math}
\usepackage{enumitem,needspace,xurl,hyperref}
\setlength{\parindent}{0pt}
\setlength{\parskip}{5pt}
\setlength{\emergencystretch}{3em}
\newcommand{\N}{\mathbb{N}}
\newcommand{\Z}{\mathbb{Z}}
\newcommand{\Q}{\mathbb{Q}}
\newcommand{\R}{\mathbb{R}}
\newcommand{\C}{\mathbb{C}}
\newcommand{\F}{\mathbb{F}}
\newcommand{\A}{\mathbb{A}}
\newcommand{\PP}{\mathbb{P}}
\newcommand{\E}{\mathbb{E}}
\newcommand{\eps}{\varepsilon}
\newcommand{\dd}{\,\mathrm{d}}
\begin{document}
\section*{UnsolvedMath ID: 2400015}
\textbf{Source problem code:} catalogue entry
\subsection{Definitions and mathematical statement}
Let $\zeta(s)$ be the meromorphic continuation of $\sum_{n\geq1}n^{-s}$ from $\operatorname{Re}s>1$, and let
\[
 \xi(s)=\tfrac12s(s-1)\pi^{-s/2}\Gamma(s/2)\zeta(s),\qquad
 \Xi(t)=\xi(\tfrac12+it).
\]
The de Bruijn--Newman family deforms this completed zeta function, with the standard constant $\Lambda$ defined by the threshold property that the deformed entire function has only real zeros precisely for deformation parameters $\lambda\geq\Lambda$.

The question is whether there are infinitely many Lehmer pairs of zeros in the sense used in that theory. A Lehmer-pair condition is a specific quantitative condition on a pair of nearby zeros, not merely the assertion that two ordinates are close. The supplied record does not provide its criterion, numerical threshold, auxiliary sums, simplicity/consecutiveness requirements, or whether a strong/weak variant is intended. Consequently it does not determine a precise set of pairs to count. The infinitude target is retained, but no criterion is invented or silently selected.

The supplied assessment expects infinitude and reports a heuristic count of order $cT$ up to height $T$, linked to Montgomery pair correlation/GUE. It specifies neither $c$, the count normalization, nor the exact pair definition, and provides no verified theorem supporting that asymptotic. It also attributes a bound written $\Lambda\leq0$ to Rodgers--Tao (2018). This direction conflicts with the standardly named nonnegativity result in the usual threshold convention; the sign and convention must be checked against a primary source rather than propagated as a verified resolution statement. Neither a Riemann-hypothesis assumption nor a resolution of infinitude is added.
\subsection{Short English statement}
Prove that the precise Lehmer-pair phenomenon occurs infinitely often among zeta zeros. The supplied summary calls this open, mentions numerical examples and GUE-based expectations, but omits the defining criterion and contains a sign/convention issue for the cited de Bruijn--Newman bound. Its no-proof-found claim is attributed to that review, not an independent exhaustive literature assessment.
\subsection{Sources}
\begin{itemize}
\item UnsolvedMath permanent ID 2400015, source code catalogue entry, ``Infinitely many Lehmer pairs.'' Registry curation and metadata: Ulam AI and UnsolvedMath Contributors, CC BY 4.0; updated 2026-08-06.
\item Wikipedia, \emph{List of unsolved problems in mathematics}, pinned revision 1366636767, analysis item 15: \url{https://en.wikipedia.org/w/index.php?title=List_of_unsolved_problems_in_mathematics&oldid=1366636767}, accessed 2026-07-29. The source adaptation is attributed under CC BY-SA 4.0. Inclusion under Unsolved problems is source status evidence only; primary-source/current-status release review remains required.
\item The supplied summary names de Bruijn, Newman, Rodgers--Tao (2018), and Montgomery pair-correlation/GUE heuristics without complete primary references. No authoritative Lehmer-pair definition or resolving citation is supplied; no missing bibliography is inferred.
\end{itemize}
\end{document}
